\section{A limit of literal Schur-remainder extraction}
\label{sec:schur-remainder}

The complete parent and a scalar expression obtained by eliminating state
coordinates have different feature requirements. This section gives an exact
example: the literal Schur remainders of one physical C2 family span an
infinite-dimensional vector space, although the complete family already has
five fixed testing features. The obstruction therefore concerns one extraction
procedure, and supplies no obstruction to a finite lift of the complete parent.

Use the complete quarter-parent $P_c$ in the regular-C2 section, with
$D_u=1+u^2$, $D_v=1+v^2$, state column
$\psi=(W,X,Y,Z)^T$, and source order $(d,n,m,C,A,D,\alpha)$.
Consider
\[
 c(t)=\left(1,\frac78+t,\frac98,\frac1{18},
                  \frac{11}8,\frac98,\frac1{18}\right),
 \qquad t\geq0,
 \qquad P_{c(t)}=\psi^TM_t(u,v)\psi.
\]
The full identity, before any state elimination, is
\begin{equation}
 P_{c(t)}=P_{c(0)}+tD_uD_vX^2.
 \label{eq:schur-parent-increment}
\end{equation}

\subsection{Positivity of the entire source family}

For $uv\ne0$, the upper $WXY$ block of $M_0$ is positive definite.
Indeed, all event packets have positive coefficients, and the restriction of
the signed $d$-packet to $Z=0$ is
\[
 \frac12\bigl[(uW+uvX-Y)^2+D_uv^2Y^2+D_vX^2\bigr].
\]
The three unweighted rows $(u,uv,-1)$, $(0,0,v)$, $(0,1,0)$ have
determinant $-uv$. Thus their positive weighted Gram matrix is positive
definite. With $U=u^2$, $V=v^2$, direct expansion gives
\[
 \det M_0=D_u^2D_v^2UV L_0(U,V),
\]
where
\begin{equation}
\begin{split}
 L_0(U,V)=\frac{121}{331776}\bigl(&880U+9317U/V+720V\\
                  &+880/V+9477V/U+720/U-21994\bigr).
\end{split}
\label{eq:seed-laurent}
\end{equation}
This is nonnegative for $U,V>0$ by weighted arithmetic--geometric mean.
The six positive integer weights sum to $21994$, and their weighted
exponent vectors sum to $(0,0)$:
\[
 880+9317-9477-720=0,
 \qquad -9317+720-880+9477=0.
\]
Consequently the weighted geometric mean of
$U,U/V,V,1/V,V/U,1/U$ is $1$. The determinant is nonnegative, and
the positive upper block makes its scalar Schur complement nonnegative.
Hence $M_0\succeq0$ for $uv\ne0$.

Continuity includes either angle equal to zero. Bihomogenizing the parent
to degree two in each angle pair includes the projective endpoints by the
same argument. Equivalently, signs are covered by the exact gauges
\[
 (u,v;W,X,Y,Z)\longmapsto
 (\epsilon u,\eta v;W,\eta X,\epsilon Y,\epsilon\eta Z),
 \qquad \epsilon,\eta\in\{-1,1\}.
\]
The additional term in \eqref{eq:schur-parent-increment} is nonnegative
for $t\geq0$, so the entire family is positive on all states and angles.
A complex state gives the sum of the values on its real and imaginary
parts; mixtures are handled by linearity. The same assertions hold over
real closed fields by their first-order polynomial formulations.

\subsection{The literal scalar remainders}

Restrict to $v=1$ and let $U=u^2>0$. Let $A_t$ denote the upper
$3\times3$ principal state block. Define
\begin{align*}
 A(U)&=26883U^2+105314U+29403,\\
 B(U)&=29000U^2+50384U+21384,\\
 D_t(U)&=A(U)+tB(U),\\
 N_t(U)&=112167(U-1)^2+t(121000U^2-189648U+81576).
\end{align*}
The determinant identities are
\[
 \det A_t=\frac{11}{20736}(U+1)D_t(U),\qquad
 \det M_t=\frac{11}{82944}(U+1)^2N_t(U).
\]
The literal Schur remainder after minimizing over $W,X,Y$ with $Z=1$
is therefore
\begin{equation}
 S_t(U)=\frac{(U+1)N_t(U)}{4D_t(U)}.
 \label{eq:literal-schur-remainder}
\end{equation}
Each coefficient of $D_t$ is positive for $t\geq0$, so these rational
functions are defined for all $U>0$.

The following exact identities control every possible cancellation:
\begin{align*}
 \operatorname{Res}_U(A,B)&=1386911105418617856,\\
 \operatorname{Res}_U(D_t,N_t)
   &=10441728(8t+11)(1000t+927)(33677t+55550)^2,\\
 D_t(-1)&=-49028,\\
 \operatorname{disc}_U(D_t)
   &=448(129472t^2+10942118t+17699275).
\end{align*}
In particular $D_t$ is coprime to $(U+1)N_t$ for every $t\geq0$.
For distinct $s,t$, a common divisor of $D_s,D_t$ would divide
$(s-t)B$ and then $A$, contradicting the first resultant.

\begin{proposition}
The functions $S_t:(0,\infty)\to\mathbb R$, $t\geq0$, are linearly
independent in the sense that every finite family with distinct parameters
is linearly independent. They cannot all be sums of squares from a single
fixed finite-dimensional space of functions of $U$.
\end{proposition}
\begin{proof}
Suppose $\sum_{i=1}^r\lambda_iS_{t_i}=0$ on $(0,\infty)$, with distinct
$t_i$. Multiplication by $4\prod_iD_{t_i}$ yields a polynomial identity.
Reduce it modulo $D_{t_i}$. Every term except the $i$th vanishes. Its
remaining product consists of $(U+1)N_{t_i}$ and the other denominators,
each coprime to $D_{t_i}$, so it is a unit in
$\mathbb R[U]/(D_{t_i})$. It follows that $\lambda_i=0$.

If $V$ is a finite-dimensional function space, the linear span of all
products of two elements of $V$ is finite dimensional, of dimension at
most $(\dim V)(\dim V+1)/2$. Every sum of squares from $V$ belongs to
that product span. The established linear independence rules out a
single such space for all $S_t$.
\end{proof}

\subsection{Why the complete parent avoids this obstruction}

The complete family already has five fixed testing features. Since $M_0$
is positive semidefinite, take the four entries of
$M_0(u,v)^{1/2}\psi$ together with $\sqrt{D_uD_v}X$. The coefficients
are $1,1,1,1,t$, respectively. This is a fixed finite feature
representation of \eqref{eq:schur-parent-increment}, including arbitrary
nonnegative source parameter $t$. The unique positive semidefinite square
root can be chosen as a semialgebraic testing function.

Thus substituting a source-dependent state minimizer into a previously
fixed feature space need not preserve that space. The literal-remainder
argument cannot be used to reject the complete-parent local construction,
nor does the five-feature example supply a global atlas for arbitrary
quantum supports. The file \nolinkurl{c2/schur_remainder_check.py}
recomputes the complete parent, determinants, gauges, resultants and
contact identities; its receipt records the precise finite checks.
